\chapter{Conclusions and future work}
The internship materials support a study of how analysis choices affect the grouping of Aβ40 descriptors. Checking the supplied script against the published method identified problems in parsing, residual-cluster handling, FES interpolation, replica pooling and weighted summaries. The modular implementation makes the revised operations testable while preserving the original sources.

The independent GMM experiment gives the main completed comparison. Removing SASA largely preserves the tested three-component partition, but agreement changes at other component counts and with initialization. The Rg--SASA correlation therefore cannot decide the feature choice by itself. The experiment concerns one fitting subset and the biased descriptor sample; broader conclusions require further sampling comparisons and valid equilibrium weighting.

The next reproduction step is to establish the FES energy units and how the replicas were combined. A bounded real-data smoke run can then test the pipeline and measure resources before larger runs retain the published parameters. Seed, feature and Wasserstein-scaling experiments should compare the same frame identities so that changes in assignments have a clear meaning.

Additional inputs would enable distinct extensions. Verified bias-offset data would support time-dependent reweighting. Atomic trajectories with a reliable frame map would allow representative structures and coordinate-based analyses using tools such as MDAnalysis and ShapeGMM \autocite{michaud2011,sasmal2024}. An external ensemble comparison would require compatible observable definitions and appropriate data access.

The report's contribution is a documented analysis that can be inspected, with a completed test of feature sensitivity and explicit limits on what remains unknown. Its code and saved outputs provide a starting point for the remaining calculations. Equilibrium conformational families and their biological interpretation must follow from those calculations, rather than from the appearance of a stable partition alone.
